% REMOVED from the checklists at the user's request (not included anywhere).
% POH supplement procedures (Low Bus Volts, Supplements 1-4, KLN 89B), verbatim, kept for reference / restoring.
%=====================================================================
\clearpage
\section{Supplements -- Emergency Procedures}
\markboth{Supplements}{}

\begin{proc}{\color{emerg}LOW BUS VOLTS WARNING (Lees Avionics Ltd Supplement No. AFMS 676002, Issue 1)}
\txt{A steady/flashing warning light is installed which will illuminate if the generator / alternator output reduces to a level where the battery supplies power to the bus-bar.}
\cond{\textbf{Before engine start}}
\chk{Check low volts warning}{ON}
\cond{\textbf{After engine start}}
\chk{Check low volts warning}{OFF}
\cond{\textbf{If warning light illuminates during flight}}
\txt{Reduce electrical load}
\txt{Battery duration approx. \textbf{30..} Mins}
\txt{Land as soon as possible.}
\pohnote{NOTE:}{Warning may illuminate with low engine rpm. Check it goes out when rpm increased.}
\end{proc}

\begin{proc}{\color{emerg}SUPPLEMENT NO. 1 -- PIPER ELECTRIC PITCH TRIM: SECTION 3 - EMERGENCY PROCEDURES}
\txt{(a) In case of malfunction, disconnect electric pitch trim by activating pitch trim switch on pilot's control yoke to OFF position.}
\txt{(b) In an emergency, electric pitch trim may be overpowered using manual pitch trim.}
\txt{(c) In cruise configuration, malfunction results in 10\textdegree\ pitch change and 200 ft altitude variation.}
\txt{(d) In approach configuration, a malfunction can result in a 5\textdegree\ pitch change and 50 ft altitude loss.}
\end{proc}

\begin{proc}{\color{emerg}SUPPLEMENT NO. 2 -- AIR CONDITIONING INSTALLATION: SECTION 3 - EMERGENCY PROCEDURES}
\txt{No change.}
\end{proc}

\begin{proc}{\color{emerg}SUPPLEMENT NO. 3 -- KING KAP 100 SERIES FLIGHT CONTROL SYSTEM: SECTION 3 - EMERGENCY PROCEDURES}
\cond{(a) SYSTEM WITH AUTOPILOT ONLY}
\txt{\quad(1) In case of Autopilot malfunction: (accomplish items a. and b. simultaneously)}
\txt{\qquad a. Airplane Control Wheel - GRASP FIRMLY and regain aircraft control.}
\txt{\qquad b. AP ENG Button - PRESS to disengage autopilot.}
\cond{(b) SYSTEMS WITH AUTOPILOT AND OPTIONAL MANUAL ELECTRIC TRIM}
\txt{\quad(1) In case of Autopilot malfunction: (accomplish items a. and b. simultaneously)}
\txt{\qquad a. Airplane Control Wheel - GRASP FIRMLY and regain aircraft control.}
\txt{\qquad b. AP DISC/TRIM INTER Switch - PRESS.}
\txt{\quad(2) In case of Manual Electric Trim malfunction:}
\txt{\qquad a. AP DISC/TRIM INTER Switch - PRESS and HOLD.}
\txt{\qquad b. PITCH TRIM Circuit Breaker - PULL.}
\txt{\qquad c. Aircraft - RETRIM manually.}
\end{proc}

\begin{proc}{\color{emerg}SUPPLEMENT NO. 4 -- AUXILIARY VACUUM SYSTEM: SECTION 3 - EMERGENCY PROCEDURES}
\subsubsection*{LOSS OF VACUUM SUCTION}
\cond{(a) Low VAC warning and VAC OFF lights illuminated:}
\chk{Vacuum suction gauge}{check to verify inop pump}
\cond{(b) If vacuum suction gauge below 4.5 inches Hg:}
\chk{Auxiliary vacuum switch}{push AUX ON}
\txt{Verify vacuum suction 4.8 to 5.2 inches Hg. and Vac and VAC OFF annunciator lights out and AUX ON light illuminated.}
\pohnote{\textit{CAUTION}}{Compass error may exceed 10\textdegree\ when auxiliary vacuum system is in operation.}
\txt{(c) Monitor electrical load - verify alternator capacity is not being exceeded as indicated by the ammeter. If required turn off non-essential electrical equipment.}
\txt{(d) Land at the earliest opportunity to have primary system repaired.}
\end{proc}

\begin{proc}{\color{emerg}BENDIX/KING KLN 89B SUPPLEMENT: SECTION III - EMERGENCY PROCEDURES / ABNORMAL PROCEDURES}
\txt{A\quad If the KLN 89B GPS information is not available or invalid, utilize remaining operational navigation equipment as required.}
\txt{B\quad If a ``RAIM NOT AVAILABLE'' message is displayed while conducting an instrument approach, terminate the approach. Execute a missed approach if required.}
\txt{C\quad If a ``RAIM NOT AVAILABLE'' message is displayed in the en route or terminal phase of flight, continue to navigate using the KLN 89B or revert to an alternate means of navigation appropriate to the route and phase of flight. When continuing to use GPS navigation, position must be verified every 15 minutes using another IFR approved navigation system.}
\txt{D\quad Refer to the KLN 89B Pilot's Guide, Appendices B and C, for appropriate pilot actions to be accomplished in response to annunciated messages.}
\end{proc}

%=====================================================================
\clearpage
\section{Supplements -- Normal Procedures}
\markboth{Supplements -- Normal}{}

\begin{proc}{\color{normal}SUPPLEMENT NO. 1 -- PIPER ELECTRIC PITCH TRIM: SECTION 4 - NORMAL PROCEDURES}
\txt{The electric trim system may be turned ON and OFF by a switch located on the left horn of the pilot's control yoke adjacent to the trim control switch. The pitch trim may be changed when the electric trim system is turned on either by moving the manual pitch trim control wheel or by operating the trim control switch on the pilot's control yoke.}
\end{proc}

\begin{proc}{\color{normal}SUPPLEMENT NO. 2 -- AIR CONDITIONING INSTALLATION: SECTION 4 - NORMAL PROCEDURES}
\txt{Prior to takeoff, the air conditioner should be checked for proper operation as follows:}
\txt{(a) Check aircraft master switch ON.}
\txt{(b) Turn the air conditioner control switch to ON and the fan switch to one of the operating positions - the AIR COND DOOR OPEN warning light will turn on, thereby indicating proper air conditioner condenser door actuation.}
\txt{(c) Turn the air conditioner control switch to OFF - the AIR COND DOOR OPEN warning light will go out, thereby indicating the air conditioner condenser door is in the up position.}
\txt{(d) If the AIR COND DOOR OPEN light does not respond as specified above, an air conditioner system or indicator bulb malfunction is indicated and further investigation should be conducted prior to flight.}
\txt{The above operational check may be performed during flight if an in flight failure is suspected.}
\txt{The condenser door light is located on the pilot's instrument panel to the right of the tachometer. The door light illuminates when the door is open and is off when the door is closed.}
\pohnote{NOTE}{To insure maximum climb performance the air conditioner must be turned off manually before takeoff to disengage the compressor and retract the condenser door. Also the air conditioner must be turned off manually before the landing approach in preparation for a possible go-around.}
\end{proc}

\begin{proc}{\color{normal}SUPPLEMENT NO. 3 -- KING KAP 100 SERIES FLIGHT CONTROL SYSTEM: SECTION 4 - NORMAL PROCEDURES}
\cond{(a) PREFLIGHT (PERFORM PRIOR TO EACH FLIGHT)}
\txt{(1) GYROS - Allow 3-4 minutes for gyros to come up to speed.}
\txt{(2) RADIO POWER Switch - ON}
\txt{(3) PREFLIGHT TEST BUTTON - PRESS momentarily and NOTE:}
\txt{\quad a. All annunciator lights on (TRIM annunciator flashing).}
\txt{\quad b. After approximately 5 seconds, all annunciator lights off except AP which will flash approximately 12 times and then remain off.}
\pohnote{NOTE}{If trim warning light stays on then the manual electric trim did not pass preflight test. The pitch trim circuit breaker should be pulled. The autopilot can still be used.}
\txt{(4) MANUAL ELECTRIC TRIM (if installed) - TEST as follows:}
\txt{\quad a. Actuate the left side of the split switch to the fore and aft positions. The trim wheel should not move on its own. Rotate the trim wheel manually against the engaged clutch to check the pilot's overpower capability.}
\txt{\quad b. Actuate right side of split switch unit to the fore and aft positions. The trim wheel should not move on its own and normal trim wheel force is required to move it manually.}
\txt{\quad c. Press the AP DISC/TRIM INTER switch down and hold. Manual Electric Trim should not operate either nose up or or nose down.}
\txt{(5) AUTOPILOT - ENGAGE by pressing AP ENG button.}
\txt{(6) CONTROL WHEEL - MOVE left and right to verify that the autopilot can be overpowered.}
\txt{(7) AP DISC/TRIM INTER Switch - PRESS. Verify that the autopilot disconnects and all modes are cancelled.}
\txt{(8) TRIM - SET to take off position.}
\cond{(b) AUTOPILOT OPERATION}
\txt{(1) Before takeoff}
\txt{\quad AP DISC/TRIM INTER Switch - PRESS.}
\txt{(2) Autopilot Engagement}
\txt{\quad AP ENG Button - PRESS. Note AP annunciator on. If no other modes are selected the autopilot will operate in the wings level mode.}
\txt{(3) Heading Changes}
\txt{\quad a. Manual Heading Changes}
\txt{\qquad 1. CWS Button - PRESS and MANEUVER aircraft to the desired heading.}
\txt{\qquad 2. CWS Button - RELEASE. Autopilot will maintain aircraft in wings level attitude.}
\pohnote{NOTE}{Aircraft heading may change in the wings level mode due to an aircraft out of trim condition.}
\txt{\quad b. Heading Hold}
\txt{\qquad 1. Heading Selector Knob - SET BUG to desired heading.}
\txt{\qquad 2. HDG Mode Selector Button - PRESS. Note HDG mode annunciator ON. Autopilot will automatically turn the aircraft to the selected heading.}
\txt{\quad c. Command Turns (Heading Hold Mode ON)}
\txt{\qquad HEADING Selector Knob - MOVE BUG to the desired heading. Autopilot will automatically turn the aircraft to the new selected heading.}
\txt{(4) NAV Coupling}
\txt{\quad a. When equipped with HSI.}
\txt{\qquad 1. Course Bearing Pointer - SET to desired course.}
\pohnote{NOTE}{When equipped with NAV 1/NAV 2 switching and NAV 2 is selected, set OBS to the desired course.}
\txt{\qquad 2. HEADING SELECTOR KNOB - SET BUG to provide desired intercept angle.}
\txt{\qquad 3. NAV Mode Selector Button - PRESS.}
\txt{\qquad If the Course Deviation Bar is greater than 2 to 3 dots: the aircraft will continue in HDG mode (or wings level if HDG not selected) with the NAV annunciator flashing; when the computed capture point is reached the HDG will disengage, the NAV annunciator will illuminate steady and the selected course will be automatically captured and tracked.}
\txt{\qquad If the D-Bar is less than 2 to 3 dots: the HDG mode will disengage upon selecting NAV mode; the NAV annunciator will illuminate steady and the capture/track sequence will automatically begin.}
\txt{\quad b. When equipped with DG}
\txt{\qquad 1. OBS Knob - SELECT desired course.}
\txt{\qquad 2. NAV Mode Selector Button - PRESS.}
\txt{\qquad 3. Heading Selector Knob - ROTATE BUG to agree with OBS course.}
\pohnote{NOTE}{When NAV is selected, the lateral operating mode will change from HDG (if selected) to wings level for 5 seconds. A 45\textdegree\ intercept angle will then be automatically established based on the position of the bug.}
\txt{\qquad If the D-Bar is greater than 2 to 3 dots: the autopilot will annunciate HDG mode (unless HDG not selected) and NAV flashing; when the computed capture point is reached the HDG annunciator will go out, the NAV annunciator will illuminate steady and the selected course will be automatically captured and tracked.}
\txt{\qquad If the D-Bar is less than 2 to 3 dots: the HDG mode will disengage upon selecting NAV mode; the NAV annunciator will illuminate steady and the capture/track sequence will automatically begin.}
\txt{(5) Approach (APR) Coupling}
\txt{\quad a. When equipped with HSI}
\txt{\qquad 1. Course Bearing Pointer - SET to desired course.}
\pohnote{NOTE}{When equipped with NAV 1/NAV 2 switching and NAV 2 is selected, set OBS to the desired course.}
\txt{\qquad 2. HEADING Selector Knob - SET BUG to provide desired intercept angle.}
\txt{\qquad 3. APR Mode Selector Button - PRESS.}
\txt{\qquad If the Course Deviation Bar is greater than 2 to 3 dots: the aircraft will continue in HDG mode (or wings level if HDG not selected) with the APR annunciator flashing; when the computed capture point is reached the HDG will disengage, the APR annunciator will illuminate steady and the selected course will be automatically captured and tracked.}
\txt{\qquad If the D-Bar is less than 2 to 3 dots: the HDG mode will disengage upon selecting APR mode; the APR annunciator will illuminate steady and the capture/track sequence will automatically begin.}
\txt{\quad b. When equipped with DG}
\txt{\qquad 1. OBS Knob - SELECT desired approach course.}
\txt{\qquad 2. APR Mode Selector Button - PRESS.}
\txt{\qquad 3. Heading Selector Knob - ROTATE Bug to agree with OBS course.}
\pohnote{NOTE}{When APR is selected, the lateral operating mode will change from HDG (if selected) to wings level for 5 seconds. A 45\textdegree\ intercept angle will then be automatically established based on the position of the bug.}
\txt{\qquad If the D-Bar is greater than 2 to 3 dots: the autopilot will annunciate HDG mode (unless HDG not selected) and APR flashing; when the computed capture point is reached the HDG annunciator will go out, the APR annunciator will illuminate steady and the selected course will be automatically captured and tracked.}
\txt{\qquad If the D-Bar is less than 2 to 3 dots: the HDG mode will disengage upon selecting APR mode; the APR annunciator will illuminate steady and the capture/track sequence will automatically begin.}
\txt{(6) BC Approach Coupling}
\txt{\quad a. When equipped with HSI}
\txt{\qquad 1. Course Bearing Pointer - SET to the ILS front course inbound heading.}
\pohnote{NOTE}{When equipped with NAV 1/NAV 2 switching and NAV 2 is selected, set OBS to the ILS front course inbound heading.}
\txt{\qquad 2. HEADING Selector Knob - SET BUG to provide desired intercept angle.}
\txt{\qquad 3. BC Mode Selector Button - PRESS.}
\txt{\qquad If the Course Deviation Bar is greater than 2 to 3 dots: the aircraft will continue in HDG mode (or wings level if HDG not selected) with BC annunciated steady and APR annunciator flashing; when the computed capture point is reached the HDG will disengage, the BC and APR annunciators will illuminate steady and the selected course will be automatically captured and tracked.}
\txt{\qquad If the D-Bar is less than 2 to 3 dots: the HDG mode will disengage upon selecting BC mode; the APR BC annunciator will illuminate steady and the capture/track sequence will automatically begin.}
\txt{\quad b. When equipped with DG}
\txt{\qquad 1. OBS Knob - SELECT the ILS front course inbound heading.}
\txt{\qquad 2. BC Mode Selector Button - PRESS.}
\txt{\qquad 3. Heading Selector Knob - ROTATE Bug to the ILS front course inbound heading.}
\pohnote{NOTE}{When BC is selected, the lateral operating mode will change from HDG (if selected) to wings level for 5 seconds. A 45\textdegree\ intercept angle will then be established based on the position of the bug.}
\txt{\qquad If the D-Bar is greater than 2 to 3 dots: the autopilot will annunciate HDG and BC modes with APR flashing; when the computed capture point is reached the HDG annunciator will go out, the APR annunciator will illuminate steady and the selected course will be automatically captured and tracked.}
\txt{\qquad If the D-Bar is less than 2 to 3 dots: the HDG mode will disengage upon selecting BC mode; the APR BC annunciator will illuminate steady and the capture/track sequence will automatically begin.}
\txt{(7) Missed Approach}
\txt{\quad a. AP DISC/TRIM INTER - PRESS to disengage AP.}
\txt{\quad b. MISSED APPROACH - EXECUTE.}
\txt{\quad c. AP ENG Button - PRESS (if AP operation is desired). Note AP annunciator ON.}
\txt{(8) Before Landing}
\txt{\quad AP DISC/TRIM INTER - PRESS to disengage AP.}
\end{proc}

\begin{proc}{\color{normal}SUPPLEMENT NO. 4 -- AUXILIARY VACUUM SYSTEM: SECTION 4 - NORMAL PROCEDURES}
\cond{(a) Preflight Check.}
\txt{1. Turn on battery switch and verify VAC OFF light illuminated.}
\pohnote{NOTE}{Due to the electrical power requirement of the auxiliary vacuum pump it is suggested that the engine be operating while making the following checks.}
\txt{2. Turn on auxiliary vacuum pump and verify AUX ON light is illuminated and electrical load (approximate 15 amps) on ammeter.}
\txt{3. Turn off auxiliary vacuum pump and verify AUX ON light extinguished.}
\cond{(b) Inflight Check. --- Prior to entering instrument flight conditions:}
\txt{1. Turn off nonessential electrical equipment.}
\txt{2. Turn on auxiliary vacuum pump and verify AUX ON light illuminated and electrical load (approximately 15 amps on ammeter.}
\txt{3. Turn off auxiliary vacuum pump and verify AUX ON light extinguished and return to normal flight.}
\pohnote{NOTE}{For maximum service life, avoid continuous non-emergency operation of the auxiliary vacuum pump.}
\end{proc}

\begin{proc}{\color{normal}BENDIX/KING KLN 89B SUPPLEMENT: SECTION IV - NORMAL PROCEDURES}
\pohnote{WARNING}{Familiarity with the en route operation of the KLN 89B does not constitute proficiency in approach operations. Do not attempt approach operations in IMC prior to attaining proficiency in the use of the KLN 89B.}
\cond{A.\quad OPERATION}
\txt{Normal operating procedures are outlined in the KLN 89B GPS Pilot's Guide, P/N 006-08786-0000, dated May 1995, (or later applicable revision). A KLN 89B Quick Reference, P/N 006-08787-0000 dated 5/95 (or later applicable revision) containing an approach sequence, operating tips and approach related messages is intended for cockpit use by the KLN 89B familiar pilot when conducting instrument approaches.}
\cond{B\quad SYSTEM ANNUNCIATORS/SWITCHES/CONTROLS}
\txt{1.\quad HSI NAV presentation (\textbf{NAV/GPS}) switch annunciator - May be used to select data for presentation on the pilot's HSI; either NAV data from the number one navigation receiver or GPS data from the KLN 89B GPS. Presentation on the HSI is also required for autopilot coupling. \textbf{NAV} is green. \textbf{GPS} is blue.}
\txt{2.\quad Message (\textbf{MSG}) annunciator - Will flash to alert the pilot of a situation that requires attention. Press the MSG button on the KLN 89B GPS to view the message. (Appendix B of the KLN 89B Pilot's Guide contains a list of all of the message page messages and their meanings). \textbf{MSG} is amber}
\txt{3\quad Waypoint (\textbf{WPT}) annunciator - Prior to reaching a waypoint in the active flight plan, the KLN 89B GPS will provide navigation along a curved path segment to ensure a smooth transition between two adjacent legs in the flight plan. This feature is called turn anticipation. Approximately 20 seconds prior to the beginning of turn anticipation the WPT annunciator will flash, going solid upon initialization of the turn, and extinguishing upon turn completion. \textbf{WPT} is amber.}
\pohnote{WARNING}{Turn anticipation is automatically disabled for FAF waypoints and those used exclusively in SID/STARS where overflight is required. For waypoints shared between SID/STARS and published en route segments (requiring overflight in the SID/STARS), proper selection on the presented waypoint page is necessary to provide adequate route protection on the SID/STARS.}
\txt{4\quad HSI course control knob - Provides analog course input to the KLN 89B in \textbf{OBS} when the NAV/GPS switch/annunciator is in \textbf{GPS}. When the NAV/GPS switch annunciation is in NAV, GPS course selection in OBS mode is digital through the use of the controls and display at the KLN 89B. The HSI course control knob must also be set to provide proper course datum to the autopilot if coupled to the KLN 89B in LEG or OBS.}
\pohnote{NOTE}{Manual HSI course centering in \textbf{OBS} using the control knob can be difficult, especially at long distances. Centering the dbar can best be accomplished by pressing [D$\rightarrow$] and then manually setting the HSI pointer to the course value prescribed in the KLN 89B displayed message.}
\txt{5\quad GPS approach (\textbf{GPS APR ARM/ACTV}) switch/annunciator - Used to a) manually select or deselect approach \textbf{ARM} (or deselect approach \textbf{ACTV}) and b) annunciate the stage of approach operation either armed (ARM) or activated (ACTV). Sequential button pushes if in \textbf{ACTV} would first result in approach \textbf{ARM} and then approach arm canceled. Subsequent button pushes will cycle between the armed state (if an approach is in the flight plan) and approach arm canceled. Approach \textbf{ACTV} cannot be selected manually. \textbf{GPS APR} and \textbf{ARM} are white. \textbf{ACTV} is green.}
\txt{6\quad RMI NAV presentation switch - May be used to select data for presentation on the RMI; either NAV 1 data from the number one navigation receiver, \textbf{NAV 2} data from the number two navigation receiver or GPS data from the KLN 89B GPS.}
\cond{C.\quad PILOTS DISPLAY}
\txt{Left/right steering information is presented on the pilot's HSI as a function of the NAV/GPS switch position.}
\cond{D\quad AUTOPILOT COUPLED OPERATION}
\txt{The KLN 89B may be coupled to the autopilot by first selecting GPS on the NAV/GPS switch. Manual selection of the desired track on the pilot's HSI course pointer is required to provide course datum to the autopilot. (Frequent manual course pointer changes may be necessary, such as in the case of flying a DME arc.) The autopilot approach mode (\textbf{APR}) \textit{should be used when} conducting a coupled GPS approach.}
\pohnote{NOTE}{\underline{Select HDG mode for DME arc intercepts.} NAV or APR coupled DME arc intercepts can result in excessive overshoots (aggravated by high ground speeds and/or intercepts from \underline{inside} the arc).}
\cond{E.\quad ALTITUDE ALERT AURAL TONES}
\txt{\textbullet\quad 1000 feet prior to reaching the selected altitude - three short tones.}
\txt{\textbullet\quad Upon reaching the selected altitude - two short tones.}
\txt{\textbullet\quad Deviating above or below the selected altitude by more than the warn altitude - four short tones.}
\cond{F.\quad APPROACH MODE SEQUENCING AND RAIM PREDICTION}
\pohnote{NOTE}{The special use airspace alert will automatically be disabled prior to flying an instrument approach to reduce the potential for message congestion.}
\txt{1.\quad Prior to arrival, select a STAR if appropriate from the \textbf{APT 7} page. Select an approach and an initial approach fix (IAF) from the \textbf{APT 8} page.}
\pohnote{NOTES}{\textbullet\ Using the outer knob, select the \textbf{ACT} (Active Flight Plan Waypoints) pages. Pull the inner knob out and scroll to the destination airport, then push the inner knob in and select the \textbf{ACT 7} or \textbf{ACT 8} page.\\[2pt]
\textbullet\ To delete or replace a SID, STAR or approach, select \textbf{FPL 0} page. Place the cursor over the name of the procedure, press \textbf{ENT} to change it, or \textbf{CLR} then \textbf{ENT} to delete it.}
\txt{2.\quad En route, check for RAIM availability at the destination airport ETA on the \textbf{OTH 3} page.}
\pohnote{NOTE}{RAIM must be available at the FAF in order to fly an instrument approach. Be prepared to terminate the approach upon loss of RAIM.}
\txt{3.\quad At 30 nm from the airport:}
\txt{\qquad a.\quad Verify automatic annunciation of APR \textbf{ARM}.}
\txt{\qquad b.\quad Note automatic dbar scaling change from $\pm$ 5.0nm to $\pm$ 1.0 nm over the next 30 seconds.}
\txt{\qquad c.\quad Update the KLN 89B altimeter baro setting as required.}
\txt{\qquad d.\quad Internally the KLN 89B will transition from en route to terminal integrity monitoring.}
\txt{4.\quad Select \textbf{NAV 4} page to fly the approach procedure.}
\txt{\qquad a.\quad If receiving radar vectors, or need to fly a procedure turn or holding pattern, fly in OBS until inbound to the FAF.}
\pohnote{NOTE}{OBS navigation is TO-FROM (like a VOR) without waypoint sequencing.}
\txt{\qquad b.\quad NoPT routes including DME arc's are flown in \textbf{LEG}. \underline{\textbf{LEG} is mandatory from the FAF to the MAP.}}
\pohnote{NOTE}{\underline{Select HDG mode for DME arc intercepts.} NAV or APR coupled DME arc intercepts can result in excessive overshoots (aggravated by high ground speeds and/or intercepts from \underline{inside} the arc).}
\pohnote{WARNING}{\textbf{Flying final outbound from an off-airport vortac on an overlay approach; beware of the DME distance increasing on final approach, and the GPS distance-to-waypoint decreasing, and not matching the numbers on the approach plate!}}
\txt{5.\quad At or before 2 nm from the FAF inbound:}
\txt{\qquad a.\quad \underline{Select the FAF as the active waypoint}, if not accomplished already.}
\txt{\qquad b.\quad \underline{Select LEG operation}.}
\txt{6.\quad Approaching the FAF inbound (within 2 nm.):}
\txt{\qquad a.\quad Verify APR \textbf{ACTV}.}
\txt{\qquad b.\quad Note automatic dbar scaling change from $\pm$ 1.0 nm to $\pm$ 0.3 nm over the 2 nm inbound to the FAF.}
\txt{\qquad c.\quad Internally the KLN 89B will transition from terminal to approach integrity monitoring.}
\txt{7.\quad Crossing the FAF and APR \textbf{ACTV} is \underline{not} annunciated:}
\txt{\qquad a.\quad \underline{Do not descend}.}
\txt{\qquad b.\quad Execute the missed approach.}
\txt{8.\quad Missed Approach:}
\txt{\qquad a.\quad \underline{Climb}}
\txt{\qquad b.\quad Navigate to the MAP (in APR \textbf{ARM} if APR \textbf{ACTV} is not available).}
\pohnote{NOTE}{There is no automatic \textbf{LEG} sequencing at the MAP.}
\txt{\qquad c.\quad After climbing in accordance with the published missed approach procedure, press [D$\rightarrow$], verify or change the desired holding fix and press \textbf{ENT}.}
\cond{GENERAL NOTES}
\txt{\textbullet\quad The data base must be up to date for instrument approach operation.}
\txt{\textbullet\quad Only \underline{one} approach can be in the flight plan at a time.}
\txt{\textbullet\quad If the destination airport is the active waypoint at the time of the instrument approach selection, the active waypoint will shift automatically to the chosen IAF.}
\txt{\textbullet\quad Checking RAIM prediction for your approach while en route using the \textbf{OTH 3} page is recommended. A self check occurs automatically within 2nm of the FAF. APR \textbf{ACTV} is inhibited without RAIM.}
\txt{\textbullet\quad Data cannot be altered, added to or deleted from the approach procedures contained in the data base. (DME arc intercepts may be relocated along the arc through the \textbf{NAV 4} or the \textbf{FPL 0} pages).}
\txt{\textbullet\quad \underline{Some approach waypoints do not appear on the approach plates (including in some instances the FAF)!}}
\txt{\textbullet\quad Waypoint suffixes in the flight plan:\\
\hspace*{2em}i - IAF\\ \hspace*{2em}f - FAF\\ \hspace*{2em}m - MAP\\ \hspace*{2em}h - missed approach holding fix.}
\txt{\textbullet\quad The DME arc IAF (arc intercept waypoint) will be a) on your present position radial off the arc VOR when you load the IAF into the flight plan, or b) the beginning of the arc if currently on a radial beyond the arc limit. To adjust the arc intercept to be compatible with a current radar vector, bring up the arc IAF waypoint in the \textbf{NAV 4} page scanning field or under the cursor on the \textbf{FPL 0} page, press \textbf{CLR}, then \textbf{ENT}. Fly the arc in LEG. adjust the HSI or CDI course pointer with reference to the desired track value on the \textbf{NAV 4} page (it will flash to remind you). Left/right dbar information is relative to the arc. Displayed distance is \underline{not} along the arc but direct to the active waypoint. (The \textbf{ARC} radial is also displayed in the lower right corner of the \textbf{NAV 4} page.)}
\txt{\textbullet\quad The DME arc IAF identifier may be unfamiliar. Example: D098G where 098 stands for the 098\textdegree\ radial off the referenced VOR, and G is the seventh letter in the alphabet indicating a 7 DME arc.}
\txt{\textbullet\quad APR \textbf{ARM} to APR \textbf{ACTV} is automatic provided:\\
\hspace*{2em}a.\ You are in APR \textbf{ARM} (normally automatic).\\
\hspace*{2em}b.\ \underline{You are in \textbf{LEG} mode}!\\
\hspace*{2em}c.\ \underline{The FAF is the active waypoint}!\\
\hspace*{2em}d.\ Within 2 n.m. of the FAF.\\
\hspace*{2em}e.\ Outside of the FAF.\\
\hspace*{2em}f.\ Inbound to the FAF.\\
\hspace*{2em}g.\ RAIM is available.}
\txt{\textbullet\quad Direct-To operation between the FAF and MAP cancels APR \textbf{ACTV}. Fly the missed approach in APR \textbf{ARM}.}
\txt{\textbullet\quad Flagged navigation inside the FAF may usually be restored (not guaranteed) by pressing the GPS APR button changing from \textbf{ACTV} to \textbf{ARM}. Fly the missed approach.}
\txt{\textbullet\quad The instrument approach using the KLN 89B may be essentially automatic starting 30 nm out (with a manual baro setting update) or it may require judicious selection of the OBS and LEG modes.}
\txt{\textbullet\quad APR \textbf{ARM} may be canceled at any time by pressing the GPS APR button. (A subsequent press will reselect it.)}
\cond{SECTION V - PERFORMANCE}
\txt{No change.}
\end{proc}

